\begin{pastexam}[label={px:gw3-097-ds-03}]{關務三等資料處理 97 年 第三題}
（申論；15 分）

動態資料結構與靜態資料結構有何不同？比較其優缺點？（15 分）
\end{pastexam}

\begin{pastsoln}
\textbf{不同之處。}
\begin{itemize}[nosep]
\item \textbf{靜態資料結構}：所需的記憶體大小在程式執行前（編譯時或宣告時）就決定好，執行期間不會
改變。 典型是陣列：宣告 \texttt{int A[100];} 之後就固定 $100$ 格，元素存在連續的記憶體裡，用註標
直接算出位址。
\item \textbf{動態資料結構}：執行期間依需要向系統要記憶體（C 的 \texttt{malloc}／\texttt{free}），
大小可以隨資料多寡增減。 典型是鏈結串列、用指標連起來的樹和圖：每個節點個別配置，靠指標串起來，
在記憶體裡不必相鄰（資結書第~5.1~節）。
\end{itemize}

\textbf{優缺點比較}（以陣列和鏈結串列為代表）：
\begin{center}
\small
\begin{tabular}{p{2.6cm}p{5cm}p{5cm}}
\toprule
 & 靜態（陣列） & 動態（鏈結串列） \\
\midrule
記憶體大小 & 事先固定。 估太大浪費，估太小會滿、要重新配置搬移 & 用多少要多少，只受系統記憶體限制 \\
額外空間 & 不需要指標 & 每個節點多一個（或兩個）指標 \\
存取第 $k$ 個 & 用位址公式直接算，$O(1)$（隨機存取） & 要從頭沿指標走 $k$ 步，$O(k)$ \\
插入、刪除 & 中間插入或刪除要搬動後面的元素，$O(n)$ & 找到位置後只改幾個指標，$O(1)$ \\
配置成本 & 一次配置，很快 & 每次插入都要配置、刪除要釋放，有時間成本；忘了釋放會記憶體流失 \\
其他 & 連續存放，快取效果好；可以做二分搜尋 & 節點分散，快取效果差；不能二分搜尋 \\
\bottomrule
\end{tabular}
\end{center}

\textbf{怎麼選}：資料量事先可以估計、主要操作是依位置讀取或搜尋時，用靜態結構（陣列）；資料量變化
大、經常在中間插入刪除時，用動態結構（鏈結串列）。 兩者也可以折衷：動態陣列滿了就配置兩倍大的新
陣列搬過去，平攤每次插入仍是 $O(1)$（資結書第~3.4~節）。 例如成績表固定 $50$ 位學生、常常查第 $k$
號，用陣列；排隊叫號系統人數不定、隨時有人插隊或離開，用鏈結串列。
\end{pastsoln}
